% ============================================================================
% APPENDIX SE: LIT-STRATIFICATION - Stratification Economics Research Integration
% ============================================================================
% Category: LIT (Literature Integration)
% Version: 1.0
% Last Updated: 2026-01-17
% Dependencies: GB (LIT-BECKER), RB (LIT-BENABOU), K (LIT-FEHR)
% Language: english
% ============================================================================

% ----------------------------------------------------------------------------
% HEADER BLOCK
% ----------------------------------------------------------------------------
% Code: SE
% Category: LIT (Literature Integration)
% Title: Stratification Economics Research Integration
% Author: EBF Framework
% Status: Complete
% Priority: High
% Evidence-Base: 30 papers (see data/paper-sources.yaml)
% ----------------------------------------------------------------------------

\chapter{LIT-STRATIFICATION: Stratification Economics Research Integration}
\label{app:lit-stratification}


\begin{tcolorbox}[colback=blue!5!white, colframe=blue!75!black,
    title=Quick Reference: Appendix SE]

\textbf{Appendix:} SE (LIT)

\textbf{Core Question:} What are the key findings in this domain?

\textbf{Key Concepts:}
\begin{itemize}[nosep]
\item Framework integration
\item Parameter validation
\item Cross-references
\end{itemize}

\textbf{Cross-References:} See related appendices below
\end{tcolorbox}

% ============================================================================
% CROSS-REFERENCE MAP
% ============================================================================
\begin{tcolorbox}[colback=blue!5!white,colframe=blue!75!black,title=Cross-Reference Map]
\textbf{Outgoing References:}
\begin{itemize}
    \item CORE-WHO (AAA): Multi-level welfare hierarchy and group membership
    \item CORE-WHAT (C): Positional goods and relative status utility
    \item CORE-HOW (B): Intergroup complementarities and competition
    \item CORE-WHEN (V): Historical context and institutional legacies
    \item CORE-WHERE (BBB): Wealth gap parameters and mobility estimates
    \item CORE-AWARE (AU): Awareness of group position and discrimination
    \item CORE-READY (AV): Barriers to economic mobility
    \item CORE-STAGE (AW): Intergenerational transmission of advantage
    \item LIT-BECKER (GB): Discrimination theory comparison
    \item LIT-BENABOU (RB): Inequality, stratification, POUM
    \item LIT-FEHR (K): Social preferences and group identity
    \item DOMAIN-LABOR (AE): Labor market discrimination
    \item DOMAIN-POL (AG): Political economy of redistribution
\end{itemize}

\textbf{Incoming References:}
\begin{itemize}
    \item DOMAIN-DEV: Development economics and colonial legacies
    \item PREDICT-INEQUALITY: Wealth gap predictions
    \item METHOD-LLMMC (AN): Parameter estimation for group disparities
\end{itemize}
\end{tcolorbox}

% ============================================================================
% CHAPTER LINKAGE
% ============================================================================
\begin{tcolorbox}[colback=green!5!white,colframe=green!75!black,title=Chapter Linkage]
\textbf{Primary Chapter:} Chapter 7 (Inequality and Stratification)

\textbf{Supporting Chapters:}
\begin{itemize}
    \item Chapter 3: Group identity and social preferences
    \item Chapter 9: Institutional design and policy
    \item Chapter 14: Wealth building interventions
    \item Chapter 15: International and development applications
\end{itemize}
\end{tcolorbox}

% ============================================================================
% ABSTRACT AND QUICK REFERENCE
% ============================================================================
\begin{abstract}
This appendix integrates Stratification Economics (SE) and its global extension (GSE) into the EBF framework. Developed by William Darity Jr. and colleagues, SE provides a distinct approach to understanding intergroup inequality that foregrounds group competition over scarce resources, the role of dominant groups in shaping institutional rules, and the persistence of hierarchies across generations. Unlike Becker's taste-based discrimination model (which predicts market erosion of discrimination) or standard human capital theory (which emphasizes individual investment), SE emphasizes that groups actively maintain positional advantages through institutional mechanisms. The EBF framework incorporates SE through six axioms (SE-1 to SE-6) covering group position theory, wealth vs. income inequality, intergenerational transmission, the Three Rs framework (Recognition, Redistribution, Repair), and policy interventions like Baby Bonds. The 2026 GSE extension adds transnational mechanisms linking origin- and destination-country hierarchies.
\end{abstract}

\begin{tcolorbox}[colback=gray!5!white,colframe=gray!75!black,title=Quick Reference]
\textbf{Core Contributions:}
\begin{itemize}
    \item Group position theory: Dominant groups defend relative position
    \item Wealth vs. income: Wealth gaps exceed income gaps (10:1 vs 2:1 Black-White)
    \item Intergenerational transmission: Inheritance perpetuates hierarchy
    \item Three Rs framework: Recognition, Redistribution, Repair
    \item Baby Bonds: Universal asset-building policy
    \item GSE (2026): Global extension with transnational mechanisms
\end{itemize}

\textbf{Key Parameters:}
\begin{itemize}
    \item Racial wealth gap (US): $\sim$10:1 White-to-Black median wealth
    \item Income gap: $\sim$2:1 White-to-Black median income
    \item Intergenerational elasticity: $\beta \approx 0.4$--$0.6$ (wealth)
    \item Colorism wage penalty: 10--15\% darker skin vs. lighter
    \item Baby Bond impact: 60--70\% reduction in wealth gap (simulated)
\end{itemize}

\textbf{Evidence Tier:} Mixed (Tier 2-5)
\begin{itemize}
    \item Wealth gap measurement: Tier 2 (survey data, administrative)
    \item Discrimination evidence: Tier 1-2 (audit studies, experiments)
    \item Policy simulations: Tier 3-4 (microsimulation)
    \item Historical mechanisms: Tier 3-5 (case studies, theoretical)
\end{itemize}
\end{tcolorbox}

% ============================================================================
% FUNDAMENTAL QUESTION
% ============================================================================
% -----------------------------------------------------------------------------
% APPENDIX SCOPE BOX
% -----------------------------------------------------------------------------

\begin{tcolorbox}[
    colback=orange!5!white,
    colframe=orange!75!black,
    title={\textbf{Appendix Scope: Ziel / In-Scope / Out-of-Scope / Constraints}},
    fonttitle=\bfseries
]

\textbf{Ziel (Objective):}
\begin{itemize}[nosep]
    \item Why do intergroup inequalities persist across generations despite civil rights advances and market competition---and what policies can effectively address the accumulated disadvantages of subordinate groups?
\end{itemize}

\textbf{In-Scope:}
\begin{itemize}[nosep]
    \item Fundamental Question
    \item Key Results and Evidence
    \item EBF Integration
    \item Worked Example: Racial Wealth Gap Intervention
\end{itemize}

\textbf{Out-of-Scope (delegiert an andere Appendices):}
\begin{itemize}[nosep]
    \item REF-GLOSSARY $\rightarrow$ Appendix G
    \item REF-GLOSSARY $\rightarrow$ Appendix G
    \item LIT-BECKER $\rightarrow$ Appendix GB
\end{itemize}

\textbf{Constraints (Anwendungsgrenzen):}
\begin{itemize}[nosep]
    \item [TODO: Define when this appendix does NOT apply]
    \item [TODO: Define required prerequisites]
    \item [TODO: Define known limitations]
\end{itemize}

\textbf{Lieferobjekte (Deliverables):}
\begin{itemize}[nosep]
    \item 5 Table(s)
    \item 5 Equation(s)
    \item 3 Axiom(s)
    \item Worked Example(s)
\end{itemize}

\end{tcolorbox}


\section{Fundamental Question}
\label{sec:se-fundamental}

\begin{tcolorbox}[colback=red!5!white,colframe=red!75!black,title=Fundamental Question]
\textbf{Why do intergroup inequalities persist across generations despite civil rights advances and market competition---and what policies can effectively address the accumulated disadvantages of subordinate groups?}
\end{tcolorbox}

Stratification Economics challenges two dominant economic explanations for persistent inequality:

\begin{enumerate}
    \item \textbf{Becker's taste-based discrimination:} Predicts market competition erodes discrimination, but gaps persist
    \item \textbf{Human capital theory:} Attributes gaps to individual skill differences, ignoring structural barriers
\end{enumerate}

SE offers an alternative: \textbf{Group Position Theory}. Dominant groups actively maintain hierarchies because:
\begin{itemize}
    \item Relative position provides utility (not just absolute consumption)
    \item Groups shape institutions to preserve advantages
    \item Intergenerational wealth transfers lock in historical advantages
    \item Discrimination is a rational strategy for dominant group members
\end{itemize}

For the EBF framework, SE is essential because it:
\begin{itemize}
    \item Centers group-level analysis (CORE-WHO) rather than only individual
    \item Explains why information and incentives alone fail to close gaps
    \item Provides policy tools (Baby Bonds, reparations) for structural intervention
    \item Extends behavioral economics to intergroup contexts
\end{itemize}

% ============================================================================
% THEORETICAL FOUNDATIONS
% ============================================================================
\section{Theoretical Foundations}
\label{sec:se-theory}

\subsection{Group Position Theory}
\label{subsec:se-group-position}

Building on Blumer (1958), SE formalizes group position as a source of utility:

\begin{equation}
U_i = u(c_i) + \lambda \cdot v(r_i - r_j)
\label{eq:se-position}
\end{equation}

where $c_i$ is individual consumption, $r_i$ is own group's relative position, $r_j$ is subordinate group's position, and $\lambda > 0$ captures positional preferences.

\begin{axiom}[Group Position Maintenance]
\label{ax:se-1}
\textbf{SE-1:} Dominant group members derive utility from relative group position and will invest resources to maintain hierarchies:
\begin{equation}
\frac{\partial U_D}{\partial (r_D - r_S)} > 0 \Rightarrow \text{invest in hierarchy maintenance}
\end{equation}
This explains why dominant groups support discriminatory institutions even when economically costly (apparent ``taste for discrimination'' is actually rational positional defense).
\end{axiom}

\textbf{EBF Integration:} Axiom SE-1 connects to CORE-WHAT (C) through positional utility and to CORE-WHO (AAA) through group-level welfare. Unlike Becker's exogenous taste, SE treats discrimination as endogenous to positional competition.

\textbf{Key Distinction from Becker (GB-2):}
\begin{itemize}
    \item Becker: Discrimination is costly ``taste'' that markets erode
    \item SE: Discrimination is rational investment in group position that persists
\end{itemize}

\subsection{Wealth vs. Income Inequality}
\label{subsec:se-wealth}

SE emphasizes that wealth inequality exceeds income inequality and has distinct dynamics:

\begin{equation}
W_{t+1} = (1 + r) \cdot W_t + s \cdot Y_t + B_t - C_t
\label{eq:se-wealth}
\end{equation}

where $W$ is wealth, $r$ is return on wealth, $Y$ is income, $s$ is savings rate, $B$ is bequests received, and $C$ is consumption.

\begin{axiom}[Wealth-Income Divergence]
\label{ax:se-2}
\textbf{SE-2:} Wealth gaps exceed income gaps because:
\begin{enumerate}
    \item Wealth accumulates across generations ($B_t$ term)
    \item Returns to wealth exceed returns to labor ($r > g$)
    \item Access to wealth-building institutions differs by group
    \item Historical expropriation created initial disparities
\end{enumerate}
US data: Black-White wealth ratio $\approx$ 1:10; income ratio $\approx$ 1:2.
\end{axiom}

\textbf{EBF Integration:} Axiom SE-2 explains why income-focused interventions (education, job training) fail to close wealth gaps---the bequest term $B_t$ perpetuates historical advantages.

\subsection{Intergenerational Transmission}
\label{subsec:se-intergenerational}

SE models how advantages transmit across generations:

\begin{equation}
W_{i,t+1} = \alpha + \beta \cdot W_{i,t} + \gamma \cdot G_i + \epsilon_{i,t}
\label{eq:se-transmission}
\end{equation}

where $G_i$ is group membership and $\gamma$ captures group-specific advantages beyond individual inheritance.

\begin{axiom}[Cumulative Disadvantage]
\label{ax:se-3}
\textbf{SE-3:} Intergroup inequality compounds across generations through multiple mechanisms:
\begin{enumerate}
    \item \textbf{Direct inheritance:} Wealth transfers at death
    \item \textbf{Inter vivos transfers:} Gifts during life (education, housing down payments)
    \item \textbf{Social capital:} Networks, connections, information
    \item \textbf{Institutional access:} Legacy admissions, credit access, neighborhoods
\end{enumerate}
Even with identical $\beta$, group differences in initial conditions ($W_0$) perpetuate indefinitely.
\end{axiom}

\textbf{EBF Integration:} Axiom SE-3 connects to CORE-STAGE (AW)---the ``behavioral change journey'' for subordinate groups is constrained by structural barriers, not just individual motivation.

\subsection{The Three Rs Framework}
\label{subsec:se-three-rs}

GSE (2026) formalizes three pillars for addressing stratification:

\begin{axiom}[Recognition-Redistribution-Repair]
\label{ax:se-4}
\textbf{SE-4:} Effective policy addressing intergroup inequality requires three components:
\begin{enumerate}
    \item \textbf{Recognition:} Acknowledge historical injustices and ongoing discrimination
    \item \textbf{Redistribution:} Transfer resources to reduce wealth gaps
    \item \textbf{Repair:} Address specific harms with targeted remedies (reparations)
\end{enumerate}
Policies addressing only one pillar are insufficient; all three must work together.
\end{axiom}

\textbf{EBF Integration:} The Three Rs framework extends EBF's intervention design by adding historical and structural dimensions beyond individual behavior change.

\subsection{Baby Bonds Policy}
\label{subsec:se-baby-bonds}

Hamilton and Darity's Baby Bonds proposal provides universal asset-building:

\begin{equation}
BB_i = f(\text{family wealth percentile}_i)
\label{eq:se-baby-bonds}
\end{equation}

where bonds range from $\sim$\$500 (wealthy families) to $\sim$\$60,000 (asset-poor families), accessible at age 18 for wealth-building activities.

\begin{axiom}[Universal Asset Building]
\label{ax:se-5}
\textbf{SE-5:} Baby Bonds address wealth inequality by:
\begin{enumerate}
    \item Providing capital at young adulthood (wealth-building period)
    \item Scaling inversely with family wealth (progressive)
    \item Being universal (political sustainability, no stigma)
    \item Targeting wealth directly (not income)
\end{enumerate}
Simulations show 60--70\% reduction in racial wealth gap within one generation.
\end{axiom}

\textbf{EBF Integration:} Baby Bonds exemplify structural intervention---changing initial conditions rather than relying on behavior change from disadvantaged individuals.

\subsection{Global Stratification Economics (GSE)}
\label{subsec:se-gse}

The 2026 GSE framework \citep{PAP-aja_pena_lopez_2026_gse} extends SE internationally:

\begin{axiom}[Transnational Stratification Mechanisms]
\label{ax:se-6}
\textbf{SE-6:} Global stratification operates through:
\begin{enumerate}
    \item \textbf{Historical genesis:} Settler colonialism, slavery, caste, extractive governance
    \item \textbf{Local institutions:} Country-specific mechanisms of dominance and resistance
    \item \textbf{Transnational links:}
    \begin{itemize}
        \item Imported stratification: Origin-country hierarchies transfer to destination
        \item Lateral mobility hypothesis: Groups move laterally in global hierarchy, not up
    \end{itemize}
\end{enumerate}
\end{axiom}

\textbf{EBF Integration:} GSE extends CORE-WHEN (V) context to include historical and transnational dimensions that shape local behavior.

% ============================================================================
% KEY RESULTS AND EVIDENCE
% ============================================================================
\section{Key Results and Evidence}
\label{sec:se-results}

\subsection{Wealth Gap Evidence}

\begin{table}[htbp]
\centering
\caption{US Racial Wealth Gap: Key Statistics}
\label{tab:se-wealth-gap}
\begin{tabular}{lccc}
\toprule
\textbf{Metric} & \textbf{White} & \textbf{Black} & \textbf{Ratio} \\
\midrule
Median wealth (2019) & \$188,200 & \$24,100 & 7.8:1 \\
Mean wealth (2019) & \$983,400 & \$142,500 & 6.9:1 \\
Homeownership rate & 73\% & 42\% & 1.7:1 \\
Median income & \$76,000 & \$46,000 & 1.7:1 \\
Zero or negative wealth & 10\% & 19\% & 0.5:1 \\
\bottomrule
\end{tabular}
\end{table}

\subsection{Discrimination Persistence}

\begin{table}[htbp]
\centering
\caption{Discrimination Evidence: Selected Studies}
\label{tab:se-discrimination}
\begin{tabular}{llcc}
\toprule
\textbf{Study} & \textbf{Finding} & \textbf{SE Support} & \textbf{Tier} \\
\midrule
Bertrand \& Mullainathan (2004) & 50\% fewer callbacks, Black names & Yes & 1 \\
Pager (2003) & Criminal record + race interaction & Yes & 1 \\
Quillian et al. (2017) & Meta: No decline since 1989 & Yes (persistence) & 1 \\
Darity et al. (2006) & Colorism wage penalty 10-15\% & Yes & 2 \\
\bottomrule
\end{tabular}
\end{table}

\subsection{Baby Bonds Simulations}

\begin{table}[htbp]
\centering
\caption{Baby Bonds Impact Projections}
\label{tab:se-baby-bonds}
\begin{tabular}{lcc}
\toprule
\textbf{Outcome} & \textbf{Baseline} & \textbf{With Baby Bonds} \\
\midrule
Black-White wealth gap & 10:1 & 3:1 to 4:1 \\
Young adults with $>$\$10K & 46\% & 85\% \\
Homeownership access & Limited & Significantly expanded \\
Student debt burden & High & Reduced \\
\bottomrule
\end{tabular}
\end{table}

% ============================================================================
% EBF INTEGRATION
% ============================================================================
\section{EBF Integration}
\label{sec:se-integration}

\subsection{CORE Framework Mapping}

\begin{table}[htbp]
\centering
\caption{Stratification Economics Mapped to 10C Framework}
\label{tab:se-9c-mapping}
\begin{tabular}{lll}
\toprule
\textbf{CORE} & \textbf{SE Contribution} & \textbf{Application} \\
\midrule
WHO (AAA) & Group-level welfare & Intergroup inequality analysis \\
WHAT (C) & Positional utility & Relative status matters \\
HOW (B) & Group competition & Zero-sum intergroup dynamics \\
WHEN (V) & Historical context & Colonial legacies, slavery \\
WHERE (BBB) & Wealth gap parameters & 10:1 ratio, mobility estimates \\
AWARE (AU) & Discrimination awareness & Perceived vs. actual barriers \\
READY (AV) & Structural barriers & Beyond individual motivation \\
STAGE (AW) & Intergenerational transmission & Cumulative disadvantage \\
\bottomrule
\end{tabular}
\end{table}

\subsection{Comparison with Other Frameworks}

\begin{table}[htbp]
\centering
\caption{Stratification Economics vs. Alternative Approaches}
\label{tab:se-comparison}
\begin{tabular}{lccc}
\toprule
\textbf{Feature} & \textbf{Becker} & \textbf{Human Capital} & \textbf{SE} \\
\midrule
Unit of analysis & Individual & Individual & Group \\
Discrimination source & Exogenous taste & Not central & Positional defense \\
Market effect & Erodes discrimination & Rewards skills & Maintains hierarchy \\
Gap persistence & Puzzle & Skill differences & Expected \\
Policy focus & Remove barriers & Education/training & Wealth transfers \\
\bottomrule
\end{tabular}
\end{table}

% ============================================================================
% WORKED EXAMPLE
% ============================================================================
\section{Worked Example: Racial Wealth Gap Intervention}
\label{sec:se-example}

\begin{tcolorbox}[colback=yellow!5!white,colframe=yellow!75!black,title=Worked Example: Comparing Policy Approaches]
\textbf{Scenario:} Design policy to reduce the 10:1 Black-White wealth gap.

\textbf{Approach 1: Human Capital (Education)}
\begin{itemize}
    \item Intervention: Increase college graduation rates
    \item Mechanism: Higher income → higher savings → wealth
    \item Problem: Even with equal education, wealth gap persists (SE-2)
    \item Evidence: Black college graduates have less wealth than White high school dropouts
\end{itemize}

\textbf{Approach 2: Becker Anti-Discrimination}
\begin{itemize}
    \item Intervention: Remove discriminatory barriers, enforce civil rights
    \item Mechanism: Equal treatment → equal outcomes
    \item Problem: Discrimination persists despite laws (SE-1)
    \item Evidence: Audit studies show no decline in hiring discrimination since 1989
\end{itemize}

\textbf{Approach 3: Stratification Economics (Baby Bonds)}
\begin{itemize}
    \item Intervention: Universal wealth grants at birth, scaled by family wealth
    \item Mechanism: Direct wealth transfer → initial conditions equalization
    \item Addresses: SE-2 (wealth vs. income), SE-3 (intergenerational transmission)
    \item Projected impact: 60-70\% gap reduction in one generation
\end{itemize}

\textbf{Approach 4: Full SE Framework (Three Rs)}
\begin{enumerate}
    \item \textbf{Recognition:} Truth commission on slavery and Jim Crow
    \item \textbf{Redistribution:} Baby Bonds + progressive wealth taxation
    \item \textbf{Repair:} Targeted reparations for documented harms
\end{enumerate}

\textbf{EBF Recommendation:} Combine approaches---behavioral interventions for individual barriers, structural interventions (Baby Bonds) for wealth gaps, recognition for legitimacy.
\end{tcolorbox}

% ============================================================================
% SUMMARY
% ============================================================================
\section{Summary}
\label{sec:se-summary}

\begin{tcolorbox}[colback=blue!5!white,colframe=blue!75!black,title=Key Takeaways]
\begin{enumerate}
    \item \textbf{SE-1 (Group Position):} Dominant groups actively maintain hierarchies for positional utility
    \item \textbf{SE-2 (Wealth-Income):} Wealth gaps (10:1) exceed income gaps (2:1); different dynamics
    \item \textbf{SE-3 (Transmission):} Multiple mechanisms transmit advantage across generations
    \item \textbf{SE-4 (Three Rs):} Recognition, Redistribution, and Repair are all necessary
    \item \textbf{SE-5 (Baby Bonds):} Universal wealth grants can dramatically reduce gaps
    \item \textbf{SE-6 (GSE):} Global extension adds transnational mechanisms and historical genesis
\end{enumerate}

\textbf{Integration Principle:} SE complements behavioral economics by showing when individual-level interventions are insufficient. Structural barriers require structural solutions---changing initial conditions, not just changing behavior.
\end{tcolorbox}

% ============================================================================
% GLOSSARY
% ============================================================================
\section{Glossary}
\label{sec:se-glossary}

Key terms defined in Appendix G (REF-GLOSSARY):

\begin{itemize}
    \item \textbf{Stratification Economics:} Approach emphasizing intergroup competition over resources
    \item \textbf{Group position theory:} Dominant groups derive utility from relative position
    \item \textbf{Wealth gap:} Difference in net worth between groups (distinct from income gap)
    \item \textbf{Baby Bonds:} Universal wealth grants at birth, scaled by family wealth
    \item \textbf{Three Rs:} Recognition, Redistribution, Repair framework
    \item \textbf{Colorism:} Discrimination based on skin tone within racial groups
    \item \textbf{Lateral mobility hypothesis:} Groups maintain relative position across borders
\end{itemize}

% ============================================================================
% CRITICAL FOUNDATIONS
% ============================================================================
\section{Critical Foundations}
\label{sec:se-foundations}

\begin{enumerate}
    \item Darity, W. Jr. (2005). Stratification Economics: The Role of Intergroup Inequality. \textit{J. Econ. \& Finance}, 29(2), 144--153.
    \item Darity, W. Jr. (2022). Position and Possessions: Stratification Economics and Intergroup Inequality. \textit{JEL}, 60(2), 405--448.
    \item Darity, W. Jr., \& Mullen, A. K. (2020). \textit{From Here to Equality}. UNC Press.
    \item Hamilton, D., \& Darity, W. Jr. (2010). Can `Baby Bonds' Eliminate the Racial Wealth Gap? \textit{Rev. Black Pol. Econ.}, 37(3-4), 207--216.
    \item Oliver, M. L., \& Shapiro, T. M. (2006). \textit{Black Wealth/White Wealth}. Routledge.
    \item Blumer, H. (1958). Race Prejudice as a Sense of Group Position. \textit{Pacific Soc. Rev.}, 1(1), 3--7.
    \item Aja, A. A., Pena, A. A., \& Lopez, M. (2026). Global Stratification Economics. \textit{NBER WP 34664}.
    \item Darity, W. Jr., \& Mason, P. L. (1998). Evidence on Discrimination in Employment. \textit{JEP}, 12(2), 63--90.
    \item Zewde, N. (2020). Universal Baby Bonds Reduce Black-White Wealth Inequality. \textit{Rev. Black Pol. Econ.}, 47(1), 3--19.
    \item Hamilton, D., \& Darity, W. Jr. (2017). The Political Economy of Education and the Racial Wealth Gap. \textit{Fed. Res. St. Louis Rev.}, 99(1), 59--76.
\end{enumerate}

% ============================================================================
% OPEN ISSUES
% ============================================================================
\section{Open Issues}
\label{sec:se-open}

\begin{enumerate}
    \item \textbf{Political feasibility:} How to build coalitions for wealth redistribution policies?
    \item \textbf{Reparations design:} What is the appropriate form and magnitude?
    \item \textbf{Global application:} How do SE insights translate to non-US contexts?
    \item \textbf{Behavioral interaction:} How do structural and behavioral interventions interact?
    \item \textbf{Measurement:} How to measure ``group position'' utility empirically?
\end{enumerate}

% ============================================================================
% REFERENCES SECTION
% ============================================================================
\section{References}
\label{sec:se-references}

\nocite{bcm_master}

This appendix integrates 30 papers from the Stratification Economics literature. Full references are maintained in the master bibliography (\texttt{bibliography/bcm\_master.bib}).

\textbf{Core SE Works (30 papers):}
\begin{itemize}
    \item \textbf{Foundational:} Darity (2005, 2016, 2022), Blumer (1958), Bobo \& Hutchings (1996)
    \item \textbf{Wealth Gap:} Oliver \& Shapiro (2006), Shapiro (2004), Hamilton \& Logan (2019)
    \item \textbf{Discrimination:} Darity \& Mason (1998), Darity et al. (2006)
    \item \textbf{Baby Bonds:} Hamilton \& Darity (2010, 2012, 2017), Zewde (2020)
    \item \textbf{Reparations:} Darity \& Mullen (2020), Darity (2008), Darity \& Frank (2003)
    \item \textbf{GSE (2026):} Aja, Pena \& Lopez (2026)
\end{itemize}

See Appendix G (REF-GLOSSARY) for term definitions and Appendix GB (LIT-BECKER) for comparison with taste-based discrimination theory.

\end{document}
